\chapter{Hiperparametri}
\label{app:hiperparametri}

\section{Hiperparametri modelov}
\label{app:hiperparametri-modelov}

Tabela~\ref{tab:hiperparametri-modelov} povzema glavne hiperparametre modelov v končnih eksperimentih.
Klasični primerjalni modeli uporabljajo agregirane tabelarične značilke, nevronske primerjalne glave uporabljajo bodisi agregirane značilke bodisi zamrznjene predstavitve, grafovski modeli pa uporabljajo grafovsko predstavitev časovnih oken.

\begin{table}
    \centering
    \caption{Glavni hiperparametri modelov v končnih eksperimentih.}
    \label{tab:hiperparametri-modelov}
    \footnotesize
    \setlength{\tabcolsep}{3pt}
    \renewcommand{\arraystretch}{0.9}
    \begin{tabularx}{\textwidth}{@{}>{\raggedright\arraybackslash}p{0.21\textwidth} >{\raggedright\arraybackslash}p{0.20\textwidth} >{\raggedright\arraybackslash}X@{}}
        \toprule
        \textbf{Model oziroma družina} & \textbf{Podatkovna predstavitev} & \textbf{Glavni hiperparametri} \\
        \midrule
        Naključni model & agregirane značilke & enakomerna naključna napoved; seme 42 \\
        Večinski model & agregirane značilke & napoved večinskega razreda na učni množici \\
        \ModelName{SVM} & agregirane značilke & RBF; $C=1{,}0$; \texttt{probability=true}; standardizacija značilk \\
        \ModelName{LightGBM} & agregirane značilke & \texttt{n\_estimators=100}; \texttt{learning\_rate=0.05}; \texttt{max\_depth=5}; seme 42 \\
        \ModelName{MLP} & agregirane značilke & skrita velikost 64; največ 200 epoh; učna stopnja $10^{-3}$; \texttt{dropout=0.1}; mini-serija 128; potrpežljivost 15 \\
        \ModelName{GazeMAE}/\ModelName{MOMENT} glave & zamrznjene predstavitve & skrita velikost 128; največ 100 epoh; učna stopnja $10^{-3}$; \texttt{dropout=0.2}; mini-serija 128; potrpežljivost 15 \\
        Grafovski modeli & grafi časovnih oken & \texttt{GCNConv}; 2 plasti; \texttt{hidden\_channels=64}; pozornostno združevanje; \texttt{dropout=0.1}; učna stopnja $5 \cdot 10^{-5}$; mini-serija 256; največ 350 epoh; potrpežljivost 20 \\
        \ModelName{GCN} & grafi časovnih oken & homogen graf; brez ločenih relacij in značilk povezav \\
        \ModelName{HeteroGCN-mean} & grafi časovnih oken & ločene relacije; združevanje s povprečenjem \\
        \ModelName{HeteroGCN-MLP} & grafi časovnih oken & ločene relacije; združevanje z \ModelName{MLP} \\
        \ModelName{HeteroGCN-MLP-w} & grafi časovnih oken & kot \ModelName{HeteroGCN-MLP}, z naučenimi utežmi povezav \\
        Grafovska konstrukcija & grafi časovnih oken & $k_t=1$; $k_s=1$; $k_f=3$; standardizirane značilke povezav \\
        \bottomrule
    \end{tabularx}
\end{table}


\section{Predhodno mrežno iskanje}
\label{app:mrezno-iskanje}

Nastavitve grafovske konstrukcije smo izbrali po majhnem predhodnem mrežnem iskanju na 3-razredni valenci, glede na izvorne razrede definirane v tabeli~\ref{tab:label-mapping}.
Iskanje je spreminjalo število grafovskih plasti, velikost skritih predstavitev ter parametre konstrukcije $k_t$, $k_s$ in $k_f$.
Primerjali smo več nastavitev števila plasti in skrite razsežnosti, vendar med njimi ni bilo velikih razlik.
Ker je bil najmanjši model med najboljšimi in najstabilnejšimi, hkrati pa je računsko najvarčnejši, smo za končne poskuse izbrali $L=2$ (število skritih plasti GNN) in $d=64$ (razsežnost skritih plasti).
Pri tej arhitekturni nastavitvi smo v mrežnem iskanju spreminjali $k_t \in \{1,2,3\}$, $k_s \in \{1,2,3\}$ in $k_f \in \{1,2,3,5\}$.
Najvišji rezultat v tem pomožnem iskanju je dosegla nastavitev $k_t=3$, $k_s=1$ in $k_f=5$, vendar so bile razlike majhne.
V končnih poskusih smo uporabili redkejšo nastavitev $k_t=1$, $k_s=1$ in $k_f=3$, ker zmanjša gostoto grafa in računsko zahtevnost ter ohrani stabilno delovanje.

% Vir podatkov: GFM-for-eyetracker/results/gnn_v2_valence_grid_search_2x64_phase2/2026-05-28_16-54-08/grid_summary.csv.
\begin{table}[H]
    \centering
    \caption{Povzetek predhodnega mrežnega iskanja hiperparametrov grafovske konstrukcije.}
    \label{tab:mrezno-iskanje-gnn}
    \footnotesize
    \begin{tabularx}{\textwidth}{@{}>{\raggedright\arraybackslash}p{0.26\textwidth} >{\centering\arraybackslash}p{0.11\textwidth} >{\centering\arraybackslash}p{0.11\textwidth} >{\centering\arraybackslash}p{0.11\textwidth} >{\raggedright\arraybackslash}X@{}}
        \toprule
        Nastavitev & $k_t$ & $k_s$ & $k_f$ & Opomba \\
        \midrule
        Preiskani prostor & $\{1,2,3\}$ & $\{1,2,3\}$ & $\{1,2,3,5\}$ & $36$ uspešno izvedenih kombinacij \\
        Najvišji makro F1 & $3$ & $1$ & $5$ & makro F1 $0{,}510$, točnost $0{,}522$ \\
        Končna nastavitev & $1$ & $1$ & $3$ & redkejši graf in manjša računska zahtevnost \\
        \bottomrule
    \end{tabularx}
\end{table}


\section{Strojna in programska oprema}
\label{app:strojna-programska-oprema}

Glavni eksperimenti so bili izvedeni na računalniku z operacijskim sistemom Linux
Ubuntu, grafično kartico NVIDIA GeForce RTX 4070 z 12~GB grafičnega pomnilnika,
procesorjem AMD Ryzen 7 5700X, 32~GB sistemskega pomnilnika in pogonom Samsung
SSD 990 PRO NVMe.
Poskusi so bili poganjani v okolju Conda \texttt{gfm}.
Uporabljene so bile različice: Python 3.10.19, PyTorch 2.9.1+cu126, CUDA 12.6, PyG 2.7.0, scikit-learn 1.7.2, LightGBM 4.6.0, NumPy 1.25.2, pandas 2.3.3, Matplotlib 3.10.6 in seaborn 0.13.2.

\section{Naučene uteži povezav}
\label{app:naucene-utezi-povezav}

V modelu \ModelName{HeteroGCN-MLP-w} značilke povezave uporabimo za izračun skalarne uteži povezave.
Naj bo $\mathbf{e}_{ij}^{(r)}$ vektor značilk povezave $(v_i, v_j)$ v relaciji $r$.
Najprej izračunamo surovo oceno:
\begin{equation}
    s_{ij}^{(r)} =
    f_w^{(r)}
    \left(
    \mathbf{e}_{ij}^{(r)}
    \right),
\end{equation}
kjer je $f_w^{(r)}$ večslojni perceptron za ustrezno družino relacij.
Za časovne povezave naprej in nazaj uporabimo isti perceptron $f_w^{\text{čas}}$, za prostorske povezave $f_w^{\text{prostor}}$, za fiksacijske povezave pa $f_w^{\text{fiksacija}}$.
Vsak od teh perceptronov ima obliko
\begin{equation}
    d_e \rightarrow 6 \rightarrow 4 \rightarrow 2 \rightarrow 1,
\end{equation}
kjer je $d_e$ razsežnost vektorja značilk povezave za izbrano relacijo in množico signalov.

Surovo oceno nato omejimo s funkcijo $\tanh$:
\begin{equation}
    \tilde{w}_{ij}^{(r)} =
    \tanh
    \left(
    s_{ij}^{(r)}
    \right).
\end{equation}
Končno utež normaliziramo po vhodnih povezavah ciljne meritve $v_j$ znotraj iste relacije:
\begin{equation}
    w_{ij}^{(r)} =
    \frac{
        \tilde{w}_{ij}^{(r)}
    }{
        \sum_{k : (v_k, v_j) \in E_r}
        \left|
        \tilde{w}_{kj}^{(r)}
        \right|
        + \epsilon
    },
\end{equation}
kjer je $\epsilon=10^{-8}$ majhna konstanta za numerično stabilnost.
Ker normalizacija ohrani predznak uteži, lahko povezava prispevek sosednjega vozlišča tudi zmanjša, ne le poveča.
